\section{Stationary-Distribution Mapping and Invariant-Measure Conditions}

This phase rewrites the stationary-distribution block in fully explicit measure-theoretic form. Objects are defined before use, theorem assumptions are stated before invocation, and algebraic inequalities are expanded line by line.

\subsection{Benchmark Fidelity Check}

\paragraph{Step 0: exact benchmark stationary-measure equation and domain.}
The main manuscript writes the stationary-distribution law of motion as
\begin{equation}
\mu'_j(B)
=
\sum_{s\in\{A,B,C\}}\int_Z \mathbf 1\{g_s(z;M)=j\}\,Q_j(B\mid z,M)\,\mu_s(dz)
+n_E^j(M)G(B),
\qquad j\in\{A,B,C\},
\label{eq:p7_main_stationary_measure}
\end{equation}
for each measurable set $B\subseteq Z$, where $g_s(z;M)$ is described in the manuscript as the optimal organizational policy. Throughout this phase we fix policy regime and prices:
\[
M\in\{0,1\},
\qquad
(w,p_m)\in\mathcal W\times\mathcal P_m,
\]
where
\[
\mathcal P_m=[\underline p_m,\overline p_m]\subset\mathbb R_+,
\qquad
\mathcal W=[\underline w,\overline w]\subset\mathbb R_+,
\]
with $0\le\underline p_m\le\overline p_m<\infty$ and $0\le\underline w\le\overline w<\infty$.

State space is compact metric $(Z,\mathcal B(Z))$, and active organizational types are $\{A,B,C\}$.

\paragraph{Step 1: exact objects and the policy-representation tension.}
Phase 4 preserved the benchmark choice-region representation
\[
\mathcal Z_{sj}(M)=\{z\in Z: j \text{ solves the Bellman problem for current type }s\}.
\]
Equation \eqref{eq:p7_main_stationary_measure}, however, is written with a single-valued policy rule $g_s$. These two representations coincide automatically only when the Bellman maximizer is unique, or when one adds an explicit measurable tie-breaking convention. We therefore proceed conditionally on a selector representation that matches \eqref{eq:p7_main_stationary_measure}, and we state that convention explicitly instead of treating it as benchmark-primitive.

\paragraph{Step 2: inherited objects used in the present phase.}
\begin{enumerate}
    \item the Phase 3 argmax correspondence
    \[
    \Gamma_s(z;M):=\arg\max_{j\in\{A,B,C,\mathrm{exit}\}}W_{sj}(z;M),
    \qquad s\in\{A,B,C\};
    \]
    \item a measurable selector representation
    \[
    g_s:Z\to\{A,B,C,\mathrm{exit}\},
    \qquad s\in\{A,B,C\},
    \]
    satisfying $g_s(z;M)\in\Gamma_s(z;M)$ for every $z$;
    \item transition kernels $Q_j(\cdot\mid z,M)$ for $j\in\{A,B,C\}$;
    \item entrant masses $n_E^A(M),n_E^B(M),n_E^C(M)$ and entrant distribution $G$.
\end{enumerate}

\paragraph{Step 3: benchmark-fidelity boundary.}
The stationary-measure mapping below is faithful to the selector-based notation in the main manuscript. It is not a proof that the set-based Phase 4 transition objects automatically become selector-based partitions without additional conditions. Whenever ties occur on economically relevant sets, the selector interpretation is a supplemental convention and must be read that way.

\paragraph{Step 4: policy-induced continuation kernels.}
For each current type $s\in\{A,B,C\}$ and next active type $j\in\{A,B,C\}$, define
\begin{equation}
K_{sj}(B\mid z,M):=\mathbf 1\{g_s(z;M)=j\}\,Q_j(B\mid z,M),
\qquad B\in\mathcal B(Z).
\label{eq:p7_Ksj}
\end{equation}

\paragraph{Step 2: measure spaces and dynamic map.}
Let $\mathcal M_+(Z)$ be finite nonnegative Borel measures and $\mathcal M(Z)$ finite signed Borel measures on $Z$.
Define
\[
\mathcal M_+^3:=\mathcal M_+(Z)^3,
\qquad
\mathcal M^3:=\mathcal M(Z)^3.
\]

For $\mu=(\mu_A,\mu_B,\mu_C)\in\mathcal M_+^3$, define
\begin{equation}
(\mathcal T_M\mu)_j(B):=\sum_{s\in\{A,B,C\}}\int_Z K_{sj}(B\mid z,M)\,\mu_s(dz),
\qquad j\in\{A,B,C\}.
\label{eq:p7_T}
\end{equation}
Define entrant inflow measure vector
\begin{equation}
\nu_j(B):=n_E^j(M)\,G(B),
\qquad j\in\{A,B,C\}.
\label{eq:p7_nu}
\end{equation}
Define one-step map
\begin{equation}
\Phi_M(\mu):=\mathcal T_M\mu+\nu.
\label{eq:p7_Phi}
\end{equation}
Stationary distribution vector is a fixed point:
\begin{equation}
\mu^*=\Phi_M(\mu^*).
\label{eq:p7_fixed_point}
\end{equation}

\subsection{Assumptions}

\begin{assumption}[SD0: Measurable selector representation]
For each current type $s\in\{A,B,C\}$, there exists a measurable mapping
\[
g_s:Z\to\{A,B,C,\mathrm{exit}\}
\]
such that $g_s(z;M)\in\Gamma_s(z;M)$ for every $z\in Z$. If the Bellman maximizer is unique, SD0 is automatic. If ties occur, SD0 is a supplemental measurable-selection convention.
\end{assumption}

\begin{assumption}[SD1: Kernel measurability]
For each $s,j\in\{A,B,C\}$ and each $B\in\mathcal B(Z)$, mapping $z\mapsto K_{sj}(B\mid z,M)$ is Borel measurable.
\end{assumption}

\begin{assumption}[SD2: Finite entrant inflow]
$n_E^A(M),n_E^B(M),n_E^C(M)$ are finite and nonnegative, and $G$ is a probability measure.
\end{assumption}

\begin{assumption}[SD3: Operator-level contraction regularity]
The linear extension of $\mathcal T_M$ from \eqref{eq:p7_T} to signed-measure vectors satisfies
\begin{equation}
\|\mathcal T_M\eta\|_{\infty,\mathrm{TV}}
\le
\bar\rho\,\|\eta\|_{\infty,\mathrm{TV}}
\qquad\text{for every }\eta\in\mathcal M^3,
\label{eq:p7_operator_contraction}
\end{equation}
for some constant $\bar\rho\in[0,1)$. This is a supplemental regularity assumption on the selector-induced stationary-distribution operator. It is not automatic from the benchmark fact that each $Q_j(\cdot\mid z,M)$ is a probability kernel.
\end{assumption}

\subsection{Norm Choice and Completeness}

For signed-measure vector $\eta=(\eta_A,\eta_B,\eta_C)\in\mathcal M^3$, define max-TV norm
\begin{equation}
\|\eta\|_{\infty,\mathrm{TV}}:=\max_{s\in\{A,B,C\}}\|\eta_s\|_{\mathrm{TV}},
\qquad
\|\eta_s\|_{\mathrm{TV}}:=|\eta_s|(Z).
\label{eq:p7_norm}
\end{equation}

\begin{lemma}[Completeness]
$(\mathcal M^3,\|\cdot\|_{\infty,\mathrm{TV}})$ is Banach, and $\mathcal M_+^3$ is closed in this Banach space.
\end{lemma}

\begin{proof}
Each $\mathcal M(Z)$ is Banach under TV norm. Finite Cartesian products remain Banach under max norm. Nonnegativity is preserved under TV limits, so $\mathcal M_+^3$ is closed.
\end{proof}

\subsection{Well-Definedness of $\Phi_M$}

\begin{lemma}[Well-definedness]
Under SD0--SD2, if $\mu\in\mathcal M_+^3$, then $\Phi_M(\mu)\in\mathcal M_+^3$.
\end{lemma}

\begin{proof}
Fix $j$ and $B\in\mathcal B(Z)$. By SD1, $z\mapsto K_{sj}(B\mid z,M)$ is measurable. Also $0\le K_{sj}(B\mid z,M)\le1$, so for finite measure $\mu_s$,
\[
0\le\int_Z K_{sj}(B\mid z,M)\,\mu_s(dz)\le\mu_s(Z)<\infty.
\]
Summing over $s$ gives finite nonnegative value for $(\mathcal T_M\mu)_j(B)$. Standard kernel integration gives countable additivity in $B$, hence $(\mathcal T_M\mu)_j$ is a finite nonnegative measure.

By SD2, $\nu_j(B)=n_E^j(M)G(B)$ is finite nonnegative. Therefore $(\Phi_M(\mu))_j=(\mathcal T_M\mu)_j+\nu_j$ is finite nonnegative. This holds for all $j$, so $\Phi_M(\mu)\in\mathcal M_+^3$.
\end{proof}

\subsection{Contraction Inequality (Full Steps)}

\begin{lemma}[Contraction of incumbent operator]
Under SD0--SD3, for any $\mu,\tilde\mu\in\mathcal M_+^3$,
\begin{equation}
\|\mathcal T_M\mu-\mathcal T_M\tilde\mu\|_{\infty,\mathrm{TV}}
\le\bar\rho\,\|\mu-\tilde\mu\|_{\infty,\mathrm{TV}}.
\label{eq:p7_contraction_T}
\end{equation}
Consequently,
\begin{equation}
\|\Phi_M(\mu)-\Phi_M(\tilde\mu)\|_{\infty,\mathrm{TV}}
\le\bar\rho\,\|\mu-\tilde\mu\|_{\infty,\mathrm{TV}}.
\label{eq:p7_contraction_Phi}
\end{equation}
\end{lemma}

\begin{proof}
Let
\[
\delta:=\mu-\tilde\mu\in\mathcal M^3.
\]
Because each $K_{sj}(B\mid\cdot,M)$ is bounded and measurable, the integrals in \eqref{eq:p7_T} extend from nonnegative measures to signed measures by Jordan decomposition. Hence $\mathcal T_M\delta$ is well defined. By linearity of integration,
\[
\mathcal T_M\mu-\mathcal T_M\tilde\mu
=
\mathcal T_M(\mu-\tilde\mu)
=
\mathcal T_M\delta.
\]
Apply SD3 with $\eta=\delta$:
\[
\|\mathcal T_M\mu-\mathcal T_M\tilde\mu\|_{\infty,\mathrm{TV}}
=
\|\mathcal T_M\delta\|_{\infty,\mathrm{TV}}
\le
\bar\rho\,\|\delta\|_{\infty,\mathrm{TV}}
=
\bar\rho\,\|\mu-\tilde\mu\|_{\infty,\mathrm{TV}}.
\]
This is exactly \eqref{eq:p7_contraction_T}.

Finally,
\[
\Phi_M(\mu)-\Phi_M(\tilde\mu)
=(\mathcal T_M\mu+\nu)-(\mathcal T_M\tilde\mu+\nu)
=\mathcal T_M\mu-\mathcal T_M\tilde\mu,
\]
so \eqref{eq:p7_contraction_Phi} follows.
\end{proof}

\subsection{Existence and Uniqueness of Invariant Measure}

\begin{proposition}[Unique invariant measure vector]
Under SD0--SD3, there exists a unique $\mu^*\in\mathcal M_+^3$ such that
\[
\mu^*=\Phi_M(\mu^*).
\]
\end{proposition}

\begin{proof}
To apply Banach fixed-point theorem, verify assumptions:
\begin{enumerate}
    \item Complete metric space: by the completeness lemma, $\mathcal M_+^3$ is complete under $\|\cdot\|_{\infty,\mathrm{TV}}$.
    \item Self-map: by well-definedness lemma, $\Phi_M:\mathcal M_+^3\to\mathcal M_+^3$.
    \item Contraction: by \eqref{eq:p7_contraction_Phi}, modulus is $\bar\rho<1$.
\end{enumerate}
Hence Banach theorem gives existence and uniqueness of fixed point $\mu^*$.
\end{proof}

\begin{proposition}[Finite-mass bound]
Let $\mu^*$ be the unique fixed point. Then
\begin{equation}
\|\mu^*\|_{\infty,\mathrm{TV}}
\le\frac{\|\nu\|_{\infty,\mathrm{TV}}}{1-\bar\rho}.
\label{eq:p7_mass_bound}
\end{equation}
\end{proposition}

\begin{proof}
From fixed-point equation:
\begin{align*}
\|\mu^*\|_{\infty,\mathrm{TV}}
&=\|\mathcal T_M\mu^*+\nu\|_{\infty,\mathrm{TV}}\\
&\le\|\mathcal T_M\mu^*\|_{\infty,\mathrm{TV}}+\|\nu\|_{\infty,\mathrm{TV}}\\
&\le\bar\rho\|\mu^*\|_{\infty,\mathrm{TV}}+\|\nu\|_{\infty,\mathrm{TV}}.
\end{align*}
Move term and divide by $1-\bar\rho>0$:
\[
\|\mu^*\|_{\infty,\mathrm{TV}}
\le\frac{\|\nu\|_{\infty,\mathrm{TV}}}{1-\bar\rho}.
\]
\end{proof}

\subsection{Interface to Transition Matrix Objects}

\paragraph{Benchmark choice regions versus selector regions.}
Phase 4 defined the benchmark choice regions as
\[
\mathcal Z_{sj}(M):=\{z\in Z:j\in\Gamma_s(z;M)\},
\qquad s,j\in\{A,B,C,\mathrm{exit}\}.
\]
Under ties, these benchmark regions may overlap. The selector representation used in the present phase instead generates the auxiliary selector regions
\[
\widehat Z_{sj}(M):=\{z\in Z:g_s(z;M)=j\},
\qquad s,j\in\{A,B,C,\mathrm{exit}\}.
\]
These two objects coincide if the Bellman maximizer is unique. If ties occur, $\widehat Z_{sj}(M)$ depends on the supplemental selector convention in SD0.

\paragraph{Selector-based transition shares.}
For any current type $s$ with $\mu_s^*(Z)>0$, define the auxiliary selector-based transition share
\begin{equation}
\widehat P_{sj}(M):=\frac{\mu_s^*(\widehat Z_{sj}(M))}{\mu_s^*(Z)}.
\label{eq:p7_Phat}
\end{equation}

\begin{lemma}[Row sum for selector-based shares]
If $\mu_s^*(Z)>0$, then
\[
\sum_{j\in\{A,B,C,\mathrm{exit}\}}\widehat P_{sj}(M)=1.
\]
\end{lemma}

\begin{proof}
Because $g_s$ is single-valued, the family $\{\widehat Z_{sj}(M):j\in\{A,B,C,\mathrm{exit}\}\}$ forms a measurable partition of $Z$:
\[
Z=\bigsqcup_j \widehat Z_{sj}(M).
\]
By finite additivity on disjoint unions,
\[
\mu_s^*(Z)=\sum_j\mu_s^*(\widehat Z_{sj}(M)).
\]
Dividing by the positive denominator $\mu_s^*(Z)$ yields
\[
1=\sum_j\frac{\mu_s^*(\widehat Z_{sj}(M))}{\mu_s^*(Z)}
=\sum_j \widehat P_{sj}(M).
\]
\end{proof}

\paragraph{Mapping back to the benchmark object.}
If the Bellman maximizer is unique $\mu_s^*$-almost surely for each current type $s$, then
\[
\widehat Z_{sj}(M)=\mathcal Z_{sj}(M)
\quad\text{up to }\mu_s^*\text{-null sets},
\]
and therefore the selector-based shares in \eqref{eq:p7_Phat} coincide with the set-based benchmark transition shares from Phase 4. Without that additional condition, the selector-based row-sum result is an auxiliary representation rather than a direct theorem about overlapping benchmark choice regions.

\subsection{What did we just do, and why does it matter?}

This phase converted the qualitative idea of a stationary distribution into an explicit operator problem on finite signed measures. The central object is the measure-valued map $\Phi_M$, which takes a candidate incumbent distribution, a selector-based policy representation, and entrant inflows, and produces the next-period incumbent distribution. Once this operator is written carefully, the stationary-distribution problem becomes the fixed-point problem
\[
\mu^*=\Phi_M(\mu^*).
\]
The detailed contraction proof matters because it does more than show existence. It shows that the stationary object used later in the equilibrium block is single valued, stable, and recoverable by iteration under the stated assumptions. Just as importantly, the proof now makes explicit where uniqueness comes from: not from the bare fact that the benchmark continuation kernels are probability measures, but from the supplemental operator-level contraction regularity imposed in SD3.

The distinction between benchmark choice regions and auxiliary selector regions is equally important. Phase 4 preserved the benchmark object
\[
\mathcal Z_{sj}(M)=\{z\in Z:j\in\Gamma_s(z;M)\},
\]
which can overlap when there are ties. Phase 7 does not silently erase that fact. Instead, it states exactly when a measurable selector representation can be used to generate disjoint regions and row-stochastic transition shares. This is what allows later phases to work with a well-defined transition interface without pretending that tie-breaking was already part of the benchmark primitive definition.

\subsection*{End-of-Section Recap}

\begin{itemize}
    \item What has been established mathematically: Phase 7 defines an explicit measure-valued mapping $\Phi_M$, proves existence and uniqueness of an invariant measure under fully stated contraction conditions, and constructs selector-based transition shares with an explicit benchmark-to-selector interface.
    \item What assumptions were used: the inherited Bellman and entry objects from Phases 3--5, the market-clearing environment from Phase 6, and the supplemental selector and operator-contraction conditions SD0--SD3 stated in this phase.
    \item What remains to be derived next: the next phase must combine the unique stationary objects, the price-clearing equations, and the entry conditions into a full stationary competitive-equilibrium definition and existence framework.
\end{itemize}
